

\section{La cristallisation épistémologique : Installation des trois dogmes}
space{0.3cm}
oindent

L'histoire de l'enseignement secondaire en France, singulièrement au tournant des XIXe et XXe
siècles, ne se réduit pas à une chronique administrative de réformes horaires ou de modifications de
programmes. Elle constitue le théâtre d'une mutation intellectuelle et sociale profonde, un moment
de « cristallisation épistémologique » où se figent les représentations collectives du savoir
légitime. La période qui s'étend de l'arrivée au pouvoir des Républicains opportunistes autour de
1880 jusqu'à la grande réforme de 1902 est décisive. C'est durant cet intervalle critique que
l'enseignement des langues anciennes — latin et grec — opère une métamorphose radicale de sa
justification sociale et pédagogique.

Confronté à la montée en puissance des sciences positives, à l'essor des langues vivantes rendues
nécessaires par la mondialisation des échanges, et à la pression démocratique pour une école plus
ouverte, l'humanisme classique aurait pu disparaître ou se marginaliser. Il n'en fut rien. Au
contraire, il s'est restructuré autour de trois dogmes inébranlables qui allaient verrouiller
l'imaginaire scolaire français pour le siècle à venir : la \textbf{Gymnastique Mentale}, le
\textbf{Texte-Monument} et l'\textbf{Élitisme} (ou la Distinction).

Ces trois piliers ne sont pas de simples arguments de circonstance ; ils forment un système
cohérent, une idéologie défensive et offensive. Le latin ne sert plus à communiquer (fin de l'usage
véhiculaire), il sert à « former l'esprit » (Gymnastique). Le texte antique n'est plus un modèle à
imiter pour produire du discours (fin de la Rhétorique), il devient un objet sacré à disséquer
scientifiquement (Texte-Monument). Enfin, cet enseignement ne vise plus l'universalité réelle, mais
la sélection d'une aristocratie de l'esprit, distinguée de la masse par la gratuité de sa culture
(Élitisme).

Ce rapport se propose d'explorer de manière exhaustive la genèse, l'installation et les implications
de ces trois dogmes. En nous appuyant sur une analyse détaillée des textes officiels, des débats
parlementaires, des manifestes intellectuels et des travaux sociologiques et historiques, nous
démontrerons comment cette cristallisation a permis à la culture classique de survivre à sa propre «
inutilité » apparente, en transformant cette dernière en vertu suprême.



\subsection{Le dogme de la gymnastique mentale : l'apologie de l'effort gratuit}

Le premier pilier de cette trinité dogmatique est l'idée de la « gymnastique mentale ». Au moment où
l'utilitarisme anglo-saxon et le pragmatisme commercial menacent de dicter les contenus de
l'enseignement, les défenseurs des humanités opèrent un repli stratégique génial : ils abandonnent
le terrain de l'utilité pratique pour investir celui de la formation cognitive.



\subsection{Les racines théoriques : la psychologie des facultés et la discipline formelle}

Pour comprendre la puissance de ce dogme, il faut remonter à ses fondements psychologiques. Le XIXe
siècle est dominé par la « psychologie des facultés » (\textit{Faculty Psychology}), héritée de la
scolastique et formalisée par des penseurs comme Thomas Reid.\cite{I-1-3-ref1} Cette théorie postule
que l'esprit humain n'est pas un tout indifférencié, mais un assemblage de « facultés » distinctes
et autonomes : la mémoire, l'attention, le jugement, le raisonnement, la volonté.



\subsection{L'analogie musculaire}

L'analogie centrale de ce modèle est celle du muscle. De même qu'un exercice physique répété
(haltères, gymnastique) développe une force musculaire générique, transférable à n'importe quelle
activité physique (porter des charges, combattre, travailler), l'exercice intellectuel sur des
objets ardus développerait une « force mentale » universelle.\cite{I-1-3-ref2} C'est la théorie de
la « discipline formelle » : peu importe \textit{ce que} l'on apprend, ce qui compte est
\textit{l'effort} que cet apprentissage exige. Le contenu est contingent ; la forme de l'exercice
est essentielle.

Dans cette perspective, le latin et le grec apparaissent comme les agrès parfaits de ce gymnase
intellectuel. Leur complexité morphologique (déclinaisons, conjugaisons), leur syntaxe rigoureuse
(système des cas, emboîtement des subordonnées) et leur éloignement culturel imposent à l'esprit une
contrainte maximale. Contrairement aux langues vivantes, qui peuvent s'acquérir par simple immersion
ou imitation (méthode naturelle), les langues mortes exigent une analyse consciente, une
décomposition logique permanente. Elles sont, par excellence, l'outil de la discipline
formelle.\cite{I-1-3-ref3}



\subsection{La résistance à la nouvelle psychologie}

Vers la fin du XIXe siècle, cette psychologie des facultés commence à être contestée par la
psychologie expérimentale naissante (Wundt, James) et par les théories de Herbart, qui nient
l'existence de facultés séparées et insistent sur l'intérêt et l'association des
idées.\cite{I-1-3-ref3} Pourtant, dans le champ pédagogique français, le dogme de la gymnastique
mentale résiste, et se renforce même. Il devient un argument politique : face à la « facilité »
supposée des sciences ou des langues modernes, le latin est le garant de la « volonté » et de
l'effort. Raoul Frary, dans son pamphlet \textit{La question du latin} (1885), bien qu'il critique
certains excès, témoigne de l'omniprésence de ce discours qui fait du latin le tuteur de
l'intelligence française.\cite{I-1-3-ref5}



\subsection{Michel bréal : la rationalisation philologique de la gymnastique}

La figure de Michel Bréal est centrale pour comprendre comment ce dogme s'est modernisé pour
survivre sous la République. Linguiste éminent, professeur au Collège de France, traducteur de la
grammaire comparée de Bopp, Bréal n'est pas un conservateur aveugle. Au contraire, il est l'un des
critiques les plus acerbes de la vieille tradition pédagogique jésuite et
napoléonienne.\cite{I-1-3-ref7}



\subsection{Contre le « perroquet », pour l'analyste}

Dans ses ouvrages majeurs, notamment \textit{De l'enseignement des langues anciennes} (1891), Bréal
s'attaque à la « routine » et au psittacisme. Il dénonce l'enseignement qui fait appel à la mémoire
mécanique au détriment de l'intelligence. Pour lui, la vieille gymnastique, celle qui consistait à
apprendre par cœur des listes de mots ou à imiter des formules cicéroniennes, est
stérile.\cite{I-1-3-ref9}

Cependant, Bréal ne rejette pas l'idée de gymnastique ; il veut la \textit{scientifiser}. Il propose
de remplacer la gymnastique de la mémoire par une gymnastique de la logique et de la sémantique.
L'étude du latin doit devenir une « leçon de mots » et une « leçon de choses »
intellectuelles.\cite{I-1-3-ref11} En analysant l'évolution du sens des mots (la sémantique,
discipline qu'il fonde), en comparant les structures grammaticales, l'élève ne se contente plus de
répéter ; il réfléchit sur les catégories de la pensée.



\subsection{Une caution scientifique pour les humanités}

L'apport de Bréal est décisif : il donne aux humanités classiques une caution scientifique moderne.
Grâce à lui, le latin n'est plus une vieillerie cléricale, mais un laboratoire de linguistique
comparée adapté à l'enfance. Il valide l'idée que l'analyse grammaticale est l'exercice suprême de
la raison. Ainsi, même réformé, l'enseignement des langues anciennes reste justifié par sa capacité
à structurer l'esprit, à lui donner des habitudes de rigueur et de précision que les autres
disciplines, jugées trop observationnelles ou trop utilitaires, ne sauraient
fournir.\cite{I-1-3-ref12}



\subsection{La crise de 1911 : henri poincaré et la sacralisation mathématique du dogme}

Le dogme de la gymnastique mentale atteint son apogée polémique et sa consécration sociale lors de
la crise de 1911. La réforme de 1902, qui avait créé des filières modernes sans latin (sections C et
D) et tenté d'établir une équivalence de dignité entre les cultures scientifique et littéraire,
provoque une réaction violente des milieux conservateurs.\cite{I-1-3-ref14}



\subsection{La ligue pour la culture française}

En 1911, Jean Richepin fonde la « Ligue pour la culture française ». Cette organisation, qui
regroupe académiciens, écrivains (Barrès, Lemaître) et notables, lance une croisade contre les «
barbares » de la réforme de 1902. Leur argument central est le déclin intellectuel de la France,
causé par l'abandon du latin.\cite{I-1-3-ref17} Le manifeste de la Ligue ne défend pas le latin pour
sa beauté littéraire, mais pour sa fonction formatrice indispensable.



\subsection{L'intervention décisive d'henri poincaré}

Le coup de maître de la Ligue est d'enrôler Henri Poincaré, mathématicien de génie universellement
respecté. Dans sa brochure \textit{Les Sciences et les Humanités} (1911), Poincaré apporte une
validation inespérée au dogme littéraire.\cite{I-1-3-ref17}

Lui, le scientifique, affirme haut et fort la « primauté des humanités classiques ». Son
raisonnement repose entièrement sur le transfert de compétences (la discipline formelle). Il
soutient que pour former un bon scientifique, il ne faut pas commencer par faire des sciences, mais
par faire du latin. Pourquoi? Parce que le scientifique a besoin non seulement de l'« esprit de
géométrie » (déduction logique), mais aussi de l'« esprit de finesse » (intuition, nuance,
compréhension globale), que seules les lettres classiques peuvent développer.\cite{I-1-3-ref17}



\subsection{Le « français des garçons de café »}

Poincaré va plus loin en liant intrinsèquement la maîtrise de la langue française à celle du latin.
Il formule une attaque célèbre contre le niveau linguistique des élèves issus des filières modernes,
qualifiant leur expression de « français des garçons de café », inapte à la précision requise par
les sciences. Selon lui, seul l'exercice de la version latine, qui oblige à peser chaque mot et à
analyser chaque nuance, permet d'acquérir la rigueur nécessaire à la formulation
scientifique.\cite{I-1-3-ref17}

Cette intervention scelle le dogme. Si le plus grand mathématicien du temps déclare que le latin est
indispensable aux mathématiques, la gymnastique mentale n'est plus une hypothèse pédagogique, c'est
une vérité établie. Le latin devient le prérequis universel à toute intelligence supérieure, le
passage obligé pour quiconque prétend à l'élite, qu'elle soit littéraire ou scientifique.



\subsection{Tableau comparatif : évolution de la justification par la gymnastique mentale}

\begin{table}[h!]\small\centering\begin{tabular}{| l | l | l | l | l |}\hline\textbf{Période} & \textbf{Acteurs Clés} & \textbf{Justification Principale} & \textbf{Nature de l'Exercice} & \textbf{Objectif Visé} \\ \hline\textbf{Avant 1880} & Jésuites, Université Napoléonienne & \textbf{Rhetorique \& Mémoire} & Discours latin, vers latins, mémorisation. & Éloquence, imitation des modèles. \\ \hline\textbf{1880-1900} & Michel Bréal, Jules Ferry & \textbf{Philologie \& Logique} & Version, grammaire comparée, sémantique. & Intelligence analytique, rigueur. \\ \hline\textbf{1900-1911} & Henri Poincaré, Ligue Culture Française & \textbf{Discipline Formelle \& Transfert} & Version comme épreuve de finesse. & Formation de l'esprit scientifique et général. \\ \hline\end{tabular}\end{table}

\subsection{Le dogme du texte-monument : la pétrification philologique}

Le deuxième dogme qui se cristallise durant cette période concerne l'objet même de l'enseignement.
Si le latin sert à muscler l'esprit, sur quoi cette gymnastique doit-elle s'exercer? La réponse qui
s'impose est celle du « Texte-Monument ». On assiste au passage d'une culture de la parole et de la
réappropriation (Rhétorique) à une culture du regard et de la conservation (Philologie).



\subsection{La mort de la rhétorique et la montée du modèle allemand}

Jusqu'au milieu du XIXe siècle, l'enseignement classique était rhétorique. Le texte ancien était
considéré comme un matériau vivant, un modèle intemporel que l'élève devait imiter pour produire ses
propres discours. On apprenait le latin pour écrire et, théoriquement, parler en latin (ou du moins
dans un style latin).

La défaite de 1870 agit comme un révélateur brutal. La France se sent humiliée par la Prusse, et
l'on attribue la supériorité allemande à sa science, à son université, et particulièrement à sa
philologie (\textit{Altertumswissenschaft}). Le modèle allemand est celui d'une science de
l'Antiquité rigoureuse, historique, critique, qui traite les textes non comme des modèles de style,
mais comme des documents historiques et linguistiques.\cite{I-1-3-ref8}



\subsection{L'arrêté de 1880 : le tournant}

Sous l'influence de réformateurs admirateurs du modèle allemand (comme Bréal ou Renan),
l'enseignement secondaire français opère sa mue. L'arrêté du 2 août 1880 et les réformes qui suivent
marquent le déclin officiel des exercices rhétoriques traditionnels.\cite{I-1-3-ref23} Le discours
latin, la narration, les vers latins — exercices de production et d'imagination — sont
progressivement marginalisés, rendus optionnels ou supprimés des examens (baccalauréat).

Ils sont remplacés par des exercices d'analyse : la \textbf{version} (traduction vers le français)
et l'\textbf{explication de texte}. Ce basculement est fondamental. On ne demande plus à l'élève
d'être un auteur latin, mais d'être un lecteur savant, capable de rendre compte du sens exact d'un
texte qu'il n'a pas le droit de modifier. Le texte devient un « Monument » intouchable, qu'il faut
nettoyer, éclairer, commenter, mais surtout pas imiter servilement ou pasticher.



\subsection{La sacralisation par la version et l'explication}

Le dogme du Texte-Monument s'incarne dans la prééminence absolue de la version. André Chervel, dans
son \textit{Histoire de l'enseignement du français}, montre comment cet exercice devient la clé de
voûte du système.\cite{I-1-3-ref26}



\subsection{La version : exercice de soumission}

La version est l'exercice de l'humilité et de la précision. Face au Monument (Cicéron, Tacite),
l'élève doit s'effacer. Il doit prouver qu'il a compris chaque nuance grammaticale, chaque allusion
historique. C'est un exercice de vérité (retrouver le sens originel) qui s'oppose à la liberté
(inventer un sens). La version latine devient l'épreuve reine du Concours Général et de
l'Agrégation, le critère ultime de l'excellence scolaire.\cite{I-1-3-ref28}



\subsection{L'explication de texte : la liturgie laïque}

L'explication de texte, qui se généralise à cette époque (venant de l'École Normale Supérieure), est
la liturgie de ce nouveau culte. Le professeur, tel un prêtre, dévoile le texte mot à mot. On
analyse la grammaire, le vocabulaire, le contexte historique. Le texte est découpé, disséqué. Cette
approche « scientifique » vise à donner à l'enseignement littéraire la même rigueur que les sciences
naturelles ou historiques. Le texte est un objet d'étude extérieur, clos sur lui-même, témoin d'une
Vérité (le Beau, le Vrai) qu'il faut contempler.\cite{I-1-3-ref23}



\subsection{L'échec programmé de la méthode directe : le refus de l'oralité}

La preuve la plus flagrante de la nature idéologique de ce dogme réside dans le sort réservé à la «
Méthode Directe » pour les langues anciennes.

Au tournant du siècle (vers 1900), une révolution pédagogique secoue l'enseignement des langues
vivantes. La « Méthode Directe » (ou méthode naturelle, intuitive), popularisée par des figures
comme Maximilian Berlitz ou Lambert Sauveur, triomphe.\cite{I-1-3-ref30} Elle repose sur l'oralité,
l'immersion, le refus de la traduction et l'apprentissage « naturel » de la
langue.\cite{I-1-3-ref32} L'administration française l'impose officiellement pour l'anglais et
l'allemand par les circulaires de 1901-1902.\cite{I-1-3-ref33}



\subsection{Pourquoi pas le latin?}

Techniquement, rien n'empêchait d'appliquer cette méthode au latin. Des pédagogues comme Lambert
Sauveur aux États-Unis l'avaient fait avec succès, enseignant le latin comme une langue vivante, par
la parole et l'interaction.\cite{I-1-3-ref35} En France même, des tentatives isolées ont existé.
Pourtant, l'institution scolaire a rejeté massivement cette approche pour les langues anciennes.

Ce rejet n'est pas pédagogique, il est épistémologique et social. Rendre le latin « vivant », oral,
communicatif, c'était le faire descendre de son piédestal de Monument. Si le latin s'apprend « sans
larmes » et « sans grammaire », comme l'anglais de Berlitz, il perd sa fonction de gymnastique
mentale (basée sur la difficulté et l'analyse) et sa sacralité patrimoniale.

Françoise Waquet, dans Le latin ou l'empire d'un signe et Parler comme un livre, analyse finement ce
phénomène.\cite{I-1-3-ref36} L'oralité du latin, encore vivante dans les milieux savants et
ecclésiastiques jusqu'au XVIIIe siècle, est tuée par l'école républicaine. Le latin doit être une
langue morte, une langue de l'œil et de l'écrit, une langue de marbre. Le dogme du Texte-Monument
exige le silence de la classe, rompu seulement par la voix du maître qui commente ou de l'élève qui
traduit laborieusement. Parler latin serait une profanation vulgaire.



\subsection{Le dogme de l'élitisme : la distinction par l'inutile}

Le troisième dogme, qui sous-tend et verrouille les deux autres, est celui de l'élitisme
sociologique, ou de la « Distinction » au sens que lui donnera Pierre Bourdieu. La période 1880-1902
est marquée par une tension sociale croissante : la démocratisation de l'école primaire et la montée
des classes moyennes (petite bourgeoisie, employés, artisans) qui réclament un accès au savoir. Face
à cette poussée, l'enseignement secondaire classique se fortifie comme une citadelle de la
distinction sociale.



\subsection{La fonction sociale du latin : bourdieu et la barrière symbolique}

L'analyse rétrospective de Pierre Bourdieu dans \textit{La Distinction} et \textit{La Noblesse
d'État} éclaire crûment les enjeux de cette période.\cite{I-1-3-ref38} Pour Bourdieu, la culture
classique fonctionne comme un « capital culturel » qui permet la reproduction des élites. Le latin
est un « signe » de classe, un shibboleth qui permet aux héritiers de se reconnaître et d'exclure
les parvenus.



\subsection{La valeur de la gratuité}

Le paradoxe apparent — pourquoi enseigner une langue inutile? — est en réalité la clé du système.
C'est précisément parce que le latin est économiquement inutile qu'il est socialement rentable.
Étudier le latin pendant sept ans (de la 6e à la Terminale) exige un temps considérable, un « loisir
» (scholè) que seules les familles aisées peuvent offrir à leurs enfants. C'est un apprentissage de
luxe, « désintéressé ».

Dans une société bourgeoise obsédée par l'utile et le profit, la culture classique pose la marque
d'une supériorité morale : celle de l'homme qui peut se permettre de perdre son temps à étudier des
choses inutiles. Le dogme de l'élitisme affirme que la véritable formation de l'esprit dirigeant
doit être coupée des réalités matérielles immédiates.\cite{I-1-3-ref40}



\subsection{Jules ferry et la « distance infranchissable »}

Il est crucial de dissiper le malentendu sur Jules Ferry. S'il est le père de l'école primaire
gratuite et obligatoire, il est aussi le gardien farouche de la séparation des ordres
d'enseignement. Dans ses discours et rapports (notamment en 1880-1881), il théorise une « distance
infranchissable » entre l'enseignement primaire (destiné au peuple) et l'enseignement secondaire
(destiné à l'élite).\cite{I-1-3-ref41}



\subsection{Le primaire supérieur vs le secondaire}

Face à la demande de scolarisation prolongée, la République développe l'Enseignement Primaire
Supérieur (EPS) et reforme l'Enseignement Spécial. Ces filières sont modernes, scientifiques,
pratiques, sans latin. Elles connaissent un grand succès auprès des classes moyennes.

Pour Ferry et les républicains de gouvernement, il ne faut surtout pas confondre ces deux ordres.
L'EPS forme des cadres intermédiaires, des sous-officiers de l'économie ; le Secondaire forme
l'état-major. La réforme de 1880 vise à rapprocher l'Enseignement Spécial de la culture générale,
mais sans lui donner le prestige du latin, afin de maintenir cette hiérarchie.\cite{I-1-3-ref43}
Plus l'enseignement moderne devient performant, plus le secondaire classique doit insister sur la
valeur irremplaçable du latin pour justifier son monopole sur les carrières de prestige (droit,
médecine, haute administration).



\subsection{La guerre des filières (1880-1902) et la victoire symbolique}

La tension culmine avec la réforme de 1902, portée par Georges Leygues. Cette réforme audacieuse
tente de moderniser le lycée en créant quatre filières au second cycle (A : Latin-Grec, B : Latin-
Langues, C : Latin-Sciences, D : Sciences-Langues) et en accordant, théoriquement, la même dignité
et le même baccalauréat à toutes.\cite{I-1-3-ref14} Pour la première fois, on peut devenir bachelier
sans latin (filière D).



\subsection{La réaction de défense}

C'est cette brèche dans le monopole qui déclenche la fureur de la Ligue pour la culture française en
1911. La violence de la réaction montre que l'enjeu n'est pas pédagogique, mais social. Les
défenseurs du latin, comme Barrès ou Richepin, voient dans la filière moderne une invasion de «
primaires » au lycée, une baisse de niveau, une « américanisation » ou une « germanisation » de la
culture française.\cite{I-1-3-ref18}

Ils construisent le mythe de la décadence : sans latin, l'élite française s'effondre. Les ingénieurs
« sans latin » sont dépeints comme des techniciens brutaux, incapables de pensée haute.



\subsection{Le triomphe du dogme}

Même si la réforme de 1902 est appliquée, le dogme de l'élitisme sort vainqueur sur le plan
symbolique. La filière A (Latin-Grec) reste la voie royale, celle des meilleurs élèves, celle de la
Distinction. Les filières sans latin sont dévalorisées, considérées comme des voies de garage ou de
rattrapage. Le latin devient un « mot de passe » social : l'avoir fait classe un individu, ne pas
l'avoir fait le déclasse. Ce système de hiérarchie par la matière, instauré à cette époque,
perdurera structurellement jusqu'aux années 1960 et symboliquement bien au-delà.

La période de 1880 à 1902 constitue bien le moment fondateur d'une « cristallisation épistémologique
» majeure dans l'histoire de l'éducation française. Face aux assauts de la modernité — scientifique,
démocratique, économique — l'enseignement des langues anciennes ne s'est pas contenté de résister ;
il s'est réinventé en se dotant d'une armure idéologique impénétrable, forgée de trois dogmes
interdépendants.

1. \textbf{La Gymnastique Mentale} a transformé l'inutilité pratique du latin en utilité cognitive
supérieure, validée par la science de l'époque (psychologie des facultés) et par des autorités
incontestables (Poincaré).

2. \textbf{Le Texte-Monument} a transformé l'objet d'étude en sanctuaire patrimonial, imposant une
pédagogie de la distance, de l'analyse et de l'écrit (version, explication), et rejetant l'oralité
vivante qui aurait pu démocratiser l'accès à la langue.

3. \textbf{L'Élitisme} a transformé cette discipline scolaire en instrument de tri social, érigeant
une barrière symbolique infranchissable entre l'élite cultivée (passée par les humanités) et la
masse instruite (passée par le moderne).

Ce système a fait preuve d'une résilience extraordinaire. En liant indissolublement l'intelligence
(Gymnastique), la culture (Monument) et le pouvoir (Élitisme), il a permis aux humanités classiques
de conserver leur prééminence au cœur du système éducatif français pendant près d'un siècle
supplémentaire, définissant ce qu'était la « Culture Générale » et modelant l'esprit des élites
républicaines bien au-delà de la IIIe République. Comprendre cette cristallisation, c'est comprendre
l'ADN profond de l'école française, ses hiérarchies implicites et sa difficulté chronique à penser
une culture commune qui ne soit pas une culture de sélection.



\subsection{Tableaux de synthèse des données}



\subsection{Tableau 1 : les vecteurs pédagogiques de la cristallisation (1880-1902)}

\begin{table}[h!]\small\centering\begin{tabular}{| p{2cm} | p{2.2cm} | p{2.2cm} | p{2.2cm} |}\hline\textbf{Dogme} & \textbf{Vecteur Institutionnel} & \textbf{Vecteur Pédagogique} & \textbf{Acteurs \& Textes de Référence} \\ \hline\textbf{Gymnastique Mentale} & \textbf{Refus de l'utilitarisme} dans le secondaire. Maintien du latin obligatoire pour les élites. & \textbf{Version Latine/Grecque}. Exercices de grammaire analytique. Rejet de la simplification. & \textbf{Michel Bréal} (\textit{De l'enseignement des langues anciennes}, 1891).\cite{I-1-3-ref7} \textbf{Henri Poincaré} (\textit{Les Sciences et les Humanités}, 1911).\cite{I-1-3-ref19} \textbf{Raoul Frary} (\textit{La question du latin}, 1885).\cite{I-1-3-ref5} \\ \hline\textbf{Texte-Monument} & \textbf{Arrêté du 2 août 1880}. Suppression du discours latin et des vers latins. & \textbf{Explication de texte}. Philologie et sémantique. Rejet de la \textbf{Méthode Directe/Orale}.\cite{I-1-3-ref33} & \textbf{André Chervel} (Histoire de la discipline).\cite{I-1-3-ref26} \textbf{Lambert Sauveur} (Méthode naturelle, rejetée pour le latin).\cite{I-1-3-ref35} \textbf{Françoise Waquet} (\textit{Le latin ou l'empire d'un signe}).\cite{I-1-3-ref36} \\ \hline\textbf{Élitisme (Distinction)} & \textbf{Réforme de 1902} (Hiérarchie implicite des sections A/B/C/D). Distinction EPS / Lycée. & Le latin comme \textbf{option de distinction}. Coût temporel élevé (7 ans d'études). & \textbf{Jules Ferry} (Discours sur la \og distance infranchissable\fg{}).\cite{I-1-3-ref41} \textbf{Ligue pour la culture française} (Manifeste de 1911).\cite{I-1-3-ref17} \textbf{Pierre Bourdieu} (\textit{La Noblesse d'État}).\cite{I-1-3-ref39} \\ \hline\end{tabular}\end{table}

\subsection{Tableau 2 : chronologie des événements clés}

\begin{table}[h!]\small\centering\begin{tabular}{| p{2.3cm} | p{3.5cm} | p{3.5cm} |}\hline\textbf{Année} & \textbf{Événement} & \textbf{Signification pour les Dogmes} \\ \hline\textbf{1880} & \textbf{Arrêté du 2 août \& Réforme Ferry} & Fin de la Rhétorique (production). Début du règne du Texte-Monument (analyse/version). Séparation stricte EPS/Secondaire (Élitisme).\cite{I-1-3-ref24} \\ \hline\textbf{1885} & \textbf{Raoul Frary, \textit{La question du latin}} & Polémique sur l'utilité du latin. Cristallisation de l'argument de la \og force mentale\fg{} face aux critiques utilitaristes.\cite{I-1-3-ref5} \\ \hline\textbf{1891} & \textbf{Michel Bréal, \textit{De l'enseignement des langues anciennes}} & Théorisation de la \og nouvelle\fg{} gymnastique : sémantique, philologique, intelligente. Rejet de la routine.\cite{I-1-3-ref7} \\ \hline\textbf{1902} & \textbf{Réforme Leygues} & Création des filières modernes (sans latin). Tentative d'égalité. Déclenche la réaction de défense des \og Classiques\fg{}.\cite{I-1-3-ref14} \\ \hline\textbf{1911} & \textbf{Fondation de la Ligue pour la culture française} & Contre-attaque des élites conservatrices. Manifeste signé par Poincaré. Sacralisation de la Gymnastique Mentale comme prérequis scientifique.\cite{I-1-3-ref17} \\ \hline\end{tabular}\end{table}